\documentclass[a4paper,12pt]{article}

\usepackage[T2A]{fontenc}
\usepackage[utf8]{inputenc}
\usepackage[russian]{babel}
\usepackage{amsmath,amssymb,amsthm,mathtools}
\usepackage{helvet}
\renewcommand{\familydefault}{\sfdefault}
\usepackage{icomma}
\usepackage[a4paper,left=20mm,right=20mm,top=20mm,bottom=22mm]{geometry}
\usepackage{microtype}
\usepackage{tikz}
\usetikzlibrary{calc,arrows.meta,positioning,shapes.geometric}
\usepackage{pgfplots}
\pgfplotsset{compat=1.18}
\usepackage{tabularx,booktabs,longtable}
\usepackage{enumitem}
\usepackage{hyperref}
\usepackage{parskip}
\usepackage{fancyhdr}
\usepackage{setspace}
\usepackage{xcolor}

\definecolor{accent}{HTML}{008080}
\definecolor{darkaccent}{HTML}{004D4D}
\definecolor{textgray}{HTML}{333333}
\definecolor{softgray}{HTML}{E8EEEE}

\hypersetup{
  colorlinks=true,
  linkcolor=darkaccent,
  urlcolor=darkaccent,
  pdftitle={Чек-лист: сложные функции, монотонность и линейная функция},
  pdfauthor={Лёвин Артём Александрович}
}

\onehalfspacing
\setlength{\parindent}{0pt}
\setlist[itemize]{leftmargin=1.6em,itemsep=0.35em,topsep=0.35em}
\setlist[enumerate]{leftmargin=1.8em,itemsep=0.55em,topsep=0.45em}
\renewcommand{\arraystretch}{1.35}

\pagestyle{fancy}
\fancyhf{}
\fancyhead[L]{\small\textcolor{textgray}{ОГЭ по математике}}
\fancyhead[R]{\small\textcolor{textgray}{Марина}}
\fancyfoot[C]{\small\thepage}
\renewcommand{\headrulewidth}{0.3pt}
\setlength{\headheight}{14pt}

\newcommand{\sect}[1]{\section*{\textcolor{darkaccent}{#1}}}
\newcommand{\subsect}[1]{\subsection*{\textcolor{accent}{#1}}}
\newcommand{\key}[1]{\textbf{\textcolor{darkaccent}{#1}}}

\tikzset{
  flowstart/.style={
    ellipse,
    draw=darkaccent,
    line width=0.9pt,
    minimum width=38mm,
    minimum height=10mm,
    align=center
  },
  flowaction/.style={
    rounded corners=3pt,
    draw=accent,
    line width=0.9pt,
    minimum width=92mm,
    minimum height=12mm,
    text width=82mm,
    align=center
  },
  flowdecision/.style={
    diamond,
    aspect=2.2,
    draw=accent,
    line width=0.9pt,
    text width=46mm,
    align=center,
    inner sep=1.5pt
  },
  flowarrow/.style={
    -{Stealth[length=3mm]},
    line width=0.9pt,
    draw=textgray
  },
  flownote/.style={
    font=\small,
    text=textgray,
    align=center
  }
}

\begin{document}

\begin{titlepage}
\thispagestyle{empty}
\centering
\vspace*{28mm}

{\Large\textcolor{darkaccent}{\textbf{Чек-лист}}\par}
\vspace{6mm}

{\LARGE\textbf{Сложные функции, монотонность}\\[2mm]
\textbf{и линейная функция}\par}

\vspace{8mm}
{\large Подготовка к ОГЭ по математике\par}

\vspace{20mm}
{\large Ученица: \textbf{Марина}\par}

\vspace{4mm}
{\large Дата: \today\par}

\vfill

{\large Лёвин Артём Александрович\par}
{\normalsize эксклюзивно для Марины\par}

\vspace{18mm}

\begin{tikzpicture}[remember picture,overlay]
  \draw[accent,line width=1.1pt]
    ($(current page.south west)+(22mm,18mm)$)
    .. controls ($(current page.south west)+(62mm,29mm)$)
    and ($(current page.south east)+(-62mm,7mm)$)
    .. ($(current page.south east)+(-22mm,18mm)$);
\end{tikzpicture}

\end{titlepage}

\sect{1. Что нужно знать о функции и подстановке}

\key{Функция} записывается, например, так:
\[
f(x)=x^2-x-6.
\]

Запись в скобках указывает, что именно подставляется вместо аргумента $x$.

\subsect{Подстановка в функцию}

Если
\[
f(x)=x^2-x-6,
\]
то
\[
f(1)=1^2-1-6=-6.
\]

При подстановке выражения $2x-3$ получаем:
\[
f(2x-3)=(2x-3)^2-(2x-3)-6.
\]

\key{Ключевой приём:} в формуле родительской функции заменяйте
\emph{каждый} $x$ тем выражением, которое записано в скобках.

\subsect{Сложная функция}

Если заданы две функции $f$ и $g$, то
\[
(f\circ g)(x)=f(g(x)).
\]

Функция $f$ здесь внешняя, а $g$ --- внутренняя.

Например, если
\[
f(x)=x^3-3x+1,
\qquad
g(x)=2x-5,
\]
то
\[
f(g(x))=(2x-5)^3-3(2x-5)+1.
\]

\subsect{Обратная задача: увидеть внешнюю и внутреннюю функции}

В выражении
\[
(5x+2)^4-2(5x+2)^2+1
\]
повторяется фрагмент $5x+2$.

Поэтому удобно взять
\[
f(t)=t^4-2t^2+1,
\qquad
g(x)=5x+2,
\]
и записать исходное выражение как
\[
f(g(x)).
\]

Аналогично
\[
\left|x^2-2x-3\right|
\]
можно представить в виде
\[
f(g(x)),
\qquad
f(t)=|t|,
\qquad
g(x)=x^2-2x-3.
\]

\subsect{ОДЗ сложной функции}

Для композиции $f(g(x))$ должны одновременно выполняться два требования:
\[
x\in D_g,
\qquad
g(x)\in D_f.
\]

Например,
\[
f(x)=\sqrt{x},
\qquad
g(x)=2-x.
\]

Тогда
\[
f(g(x))=\sqrt{2-x}.
\]

По условию существования квадратного корня:
\[
2-x\ge0,
\]
откуда
\[
x\le2.
\]

Следовательно,
\[
D_{f\circ g}=(-\infty,2].
\]

\clearpage
\thispagestyle{fancy}

\begin{center}
{\Large\textcolor{darkaccent}{
\textbf{Алгоритм: как найти композицию функций}
}}
\end{center}

\vspace{8mm}

\begin{center}
\begin{tikzpicture}[node distance=10mm]

\node[flowstart] (s) {Начало};

\node[flowaction,below=of s] (read)
{Прочитайте запись $f(g(x))$: $f$ --- внешняя функция, $g$ --- внутренняя};

\node[flowaction,below=of read] (domain)
{Запишите ограничения: $x\in D_g$ и условие $g(x)\in D_f$};

\node[flowaction,below=of domain] (subst)
{В формуле $f(x)$ замените каждый $x$ выражением $g(x)$};

\node[flowaction,below=of subst] (simp)
{Аккуратно упростите полученное выражение};

\node[flowaction,below=of simp] (check)
{Проверьте ОДЗ и укажите область определения композиции};

\node[flowstart,below=of check] (e) {Готово};

\draw[flowarrow] (s)--(read);
\draw[flowarrow] (read)--(domain);
\draw[flowarrow] (domain)--(subst);
\draw[flowarrow] (subst)--(simp);
\draw[flowarrow] (simp)--(check);
\draw[flowarrow] (check)--(e);

\end{tikzpicture}
\end{center}

\clearpage

\sect{2. Монотонность функций}

\subsect{Базовые функции}

На тех промежутках, где функции определены, полезно помнить:

\begin{itemize}
  \item $y=x$ --- возрастающая функция;
  \item $y=kx+b$: при $k>0$ возрастает, при $k<0$ убывает, при $k=0$ постоянна;
  \item $y=\sqrt{x}$ возрастает на $[0,+\infty)$;
  \item $y=-\sqrt{x}$ убывает на $[0,+\infty)$;
  \item $y=\dfrac1x$ убывает на каждом из промежутков
  $(-\infty,0)$ и $(0,+\infty)$;
  \item $y=\dfrac1{x-a}$ убывает на каждом из промежутков
  $(-\infty,a)$ и $(a,+\infty)$.
\end{itemize}

\subsect{Монотонность композиции}

Пусть обе функции монотонны на тех промежутках, где рассматривается композиция.

Тогда действует таблица:

\begin{center}
\begin{tabularx}{0.86\linewidth}{@{}XXc@{}}
\toprule
Внешняя $f$ & Внутренняя $g$ & $f\circ g$ \\
\midrule
возрастающая & возрастающая & возрастающая \\
возрастающая & убывающая & убывающая \\
убывающая & возрастающая & убывающая \\
убывающая & убывающая & возрастающая \\
\bottomrule
\end{tabularx}
\end{center}

\key{Важно:} композиция двух убывающих функций является возрастающей.

\subsect{Пример 1}

Рассмотрим функцию
\[
y=\sqrt{-2x}.
\]

Разложение:
\[
f(t)=\sqrt{t}
\quad\text{--- возрастает},
\]
\[
g(x)=-2x
\quad\text{--- убывает}.
\]

Следовательно, композиция $f(g(x))$ убывает.

ОДЗ:
\[
-2x\ge0,
\]
откуда
\[
x\le0.
\]

Итак, функция убывает на
\[
(-\infty,0].
\]

\subsect{Пример 2}

Рассмотрим функцию
\[
y=-\sqrt{x-5}.
\]

Здесь
\[
f(t)=-\sqrt{t}
\quad\text{--- убывает},
\]
\[
g(x)=x-5
\quad\text{--- возрастает}.
\]

Композиция убывает.

ОДЗ:
\[
x-5\ge0,
\]
поэтому
\[
x\ge5.
\]

Следовательно, функция убывает на
\[
[5,+\infty).
\]

\subsect{Сумма монотонных функций}

Если две функции возрастают на одном и том же промежутке, их сумма тоже возрастает.

Если обе функции убывают на одном и том же промежутке, их сумма тоже убывает.

Для функций с разным характером монотонности общего вывода по этому правилу сделать нельзя.

Например,
\[
y=\frac1{x+2}-\sqrt{x}.
\]

Область определения:
\[
x\ge0.
\]

На $[0,+\infty)$ функция
\[
\frac1{x+2}
\]
убывает.

Функция
\[
-\sqrt{x}
\]
также убывает.

Следовательно,
\[
y=\frac1{x+2}-\sqrt{x}
\]
убывает на
\[
[0,+\infty).
\]

\clearpage
\thispagestyle{fancy}

\begin{center}
{\Large\textcolor{darkaccent}{
\textbf{Алгоритм: как определить монотонность}
}}
\end{center}

\vspace{5mm}

\begin{center}
\begin{tikzpicture}[node distance=9mm]

\node[flowstart] (s) {Начало};

\node[flowaction,below=of s] (odz)
{Найдите область определения и рабочий промежуток};

\node[flowdecision,below=12mm of odz] (comp)
{Видна композиция $f(g(x))$?};

\node[
  flowaction,
  minimum width=62mm,
  text width=54mm,
  below=14mm of comp,
  xshift=-38mm
] (parts)
{Определите монотонность внешней и внутренней функций};

\node[
  flowaction,
  minimum width=62mm,
  text width=54mm,
  below=14mm of comp,
  xshift=38mm
] (sum)
{Представьте выражение как сумму известных функций};

\node[
  flowaction,
  minimum width=62mm,
  text width=54mm,
  below=17mm of parts
] (rule)
{Композиция: одинаковый характер $\Rightarrow$ рост; разный $\Rightarrow$ убывание};

\node[
  flowdecision,
  text width=34mm,
  below=17mm of sum
] (same)
{Слагаемые монотонны одинаково?};

\node[
  flowaction,
  minimum width=48mm,
  text width=40mm,
  below=14mm of same,
  xshift=12mm
] (sumres)
{Сумма имеет тот же характер};

\node[
  flowaction,
  minimum width=48mm,
  text width=40mm,
  below=14mm of same,
  xshift=-38mm
] (other)
{Правила суммы недостаточно};

\node[flowstart,below=34mm of rule] (e1) {Вывод};
\node[flowstart,below=13mm of sumres] (e2) {Вывод};

\draw[flowarrow] (s)--(odz);
\draw[flowarrow] (odz)--(comp);

\draw[flowarrow]
  (comp.west) -|
  node[flownote,pos=0.25,above]{да}
  (parts.north);

\draw[flowarrow]
  (comp.east) -|
  node[flownote,pos=0.25,above]{нет}
  (sum.north);

\draw[flowarrow] (parts)--(rule);
\draw[flowarrow] (rule)--(e1);

\draw[flowarrow] (sum)--(same);

\draw[flowarrow]
  (same.south east) --
  node[flownote,right]{да}
  (sumres.north);

\draw[flowarrow]
  (same.south west) --
  ++(-5mm,-5mm) -|
  node[flownote,pos=0.2,left]{нет}
  (other.north);

\draw[flowarrow] (sumres)--(e2);

\end{tikzpicture}
\end{center}

\clearpage

\sect{3. Линейная функция $y=kx+b$}

\subsect{Коэффициенты и монотонность}

Для линейной функции
\[
y=kx+b
\]
коэффициент $k$ отвечает за наклон графика:

\[
\begin{cases}
k>0, & \text{функция возрастает},\\
k<0, & \text{функция убывает},\\
k=0, & \text{функция постоянна}.
\end{cases}
\]

Число $b$ является ординатой точки пересечения графика с осью $Oy$.

\subsect{Точки пересечения с осями}

Рассмотрим график
\[
y=2x-1.
\]

\key{Пересечение с осью $Ox$:}

На оси $Ox$
\[
y=0.
\]

Поэтому
\[
0=2x-1,
\]
откуда
\[
x=\frac12.
\]

Точка пересечения:
\[
\left(\frac12,0\right).
\]

\key{Пересечение с осью $Oy$:}

На оси $Oy$
\[
x=0.
\]

Подставляем:
\[
y=2\cdot0-1=-1.
\]

Точка пересечения:
\[
(0,-1).
\]

\begin{center}
\begin{tikzpicture}
\begin{axis}[
  width=0.72\linewidth,
  height=0.42\linewidth,
  axis lines=middle,
  xmin=-2,
  xmax=3,
  ymin=-3,
  ymax=4,
  unit vector ratio=1 1,
  samples=220,
  clip=false,
  xlabel={$x$},
  ylabel={$y$},
  xtick={-2,-1,0,0.5,1,2,3},
  ytick={-3,-2,-1,0,1,2,3,4},
  grid=both,
  major grid style={
    line width=.2pt,
    draw=softgray
  }
]

\addplot[
  accent,
  thick,
  domain=-2:3
]{2*x-1};

\addplot[
  only marks,
  mark=*,
  mark size=2.2pt
] coordinates {
  (0.5,0)
  (0,-1)
};

\node[anchor=south west]
  at (axis cs:0.5,0)
  {$\left(\frac12,0\right)$};

\node[anchor=north east]
  at (axis cs:0,-1)
  {$(0,-1)$};

\end{axis}
\end{tikzpicture}
\end{center}

\subsect{Как найти формулу прямой по двум точкам}

Пусть прямая проходит через точки
\[
A(x_1,y_1),
\qquad
B(x_2,y_2),
\qquad
x_1\ne x_2.
\]

Можно использовать систему:
\[
\begin{cases}
y_1=kx_1+b,\\
y_2=kx_2+b.
\end{cases}
\]

Также можно сразу найти
\[
k=\frac{y_2-y_1}{x_2-x_1},
\]
а затем
\[
b=y_1-kx_1.
\]

\subsect{Пример}

Пусть
\[
A(-1,4),
\qquad
B(3,-4).
\]

Находим коэффициент $k$:
\[
k=
\frac{-4-4}{3-(-1)}
=
\frac{-8}{4}
=
-2.
\]

Теперь найдём $b$:
\[
b=4-(-2)(-1)=2.
\]

Следовательно,
\[
y=-2x+2.
\]

Поскольку
\[
k=-2<0,
\]
функция убывает.

\clearpage
\thispagestyle{fancy}

\begin{center}
{\Large\textcolor{darkaccent}{
\textbf{Алгоритм: формула прямой по двум точкам}
}}
\end{center}

\vspace{8mm}

\begin{center}
\begin{tikzpicture}[node distance=10mm]

\node[flowstart] (s) {Начало};

\node[flowaction,below=of s] (pts)
{Запишите координаты $A(x_1,y_1)$ и $B(x_2,y_2)$; проверьте $x_1\ne x_2$};

\node[flowaction,below=of pts] (k)
{Найдите наклон:
$\displaystyle k=\frac{y_2-y_1}{x_2-x_1}$};

\node[flowaction,below=of k] (b)
{Найдите свободный член:
$b=y_1-kx_1$};

\node[flowaction,below=of b] (eq)
{Запишите уравнение $y=kx+b$};

\node[flowaction,below=of eq] (verify)
{Проверьте подстановкой обе исходные точки};

\node[flowstart,below=of verify] (e) {Готово};

\draw[flowarrow] (s)--(pts);
\draw[flowarrow] (pts)--(k);
\draw[flowarrow] (k)--(b);
\draw[flowarrow] (b)--(eq);
\draw[flowarrow] (eq)--(verify);
\draw[flowarrow] (verify)--(e);

\end{tikzpicture}
\end{center}

\clearpage

\sect{4. Типичные ошибки и проверка себя}

\begin{itemize}

  \item В $f(g(x))$ подставляют формулу $g(x)$ только в одно место,
  хотя заменить необходимо каждый $x$ в формуле $f(x)$.

  \item Путают порядок композиции:
  \[
  f(g(x))
  \quad\text{и}\quad
  g(f(x))
  \]
  в общем случае дают разные функции.

  \item Забывают ОДЗ:
  подкоренное выражение должно быть неотрицательным,
  знаменатель должен быть отличен от нуля.

  \item Считают, что композиция двух убывающих функций убывает.
  Верное правило:
  композиция двух убывающих функций возрастает.

  \item Для суммы функций делают вывод по одному слагаемому.
  Правило суммы применяется напрямую, когда оба слагаемых
  монотонны одинаково на одном промежутке.

  \item В точке пересечения с осью $Ox$ забывают условие
  \[
  y=0.
  \]

  \item В точке пересечения с осью $Oy$ забывают условие
  \[
  x=0.
  \]

  \item При восстановлении прямой по двум точкам находят $k$,
  но забывают вычислить $b$.

\end{itemize}

\subsect{Мини-памятка}

\begin{center}
\begin{tabularx}{\linewidth}{@{}lX@{}}
\toprule
Ситуация & Что делать \\
\midrule

$f(g(x))$
&
Внешняя функция --- $f$.
Подставить $g(x)$ вместо каждого $x$ в формуле $f$.
\\

Корень
&
Требовать, чтобы подкоренное выражение было $\ge0$.
\\

Дробь
&
Требовать, чтобы знаменатель был $\ne0$.
\\

Пересечение с $Ox$
&
Положить $y=0$.
\\

Пересечение с $Oy$
&
Положить $x=0$.
\\

Прямая через две точки
&
Найти $k$, затем $b$, записать $y=kx+b$ и проверить обе точки.
\\

\bottomrule
\end{tabularx}
\end{center}

\clearpage

\sect{5. Тренировочный блок}

\textbf{Средний уровень}

\begin{enumerate}

  \item
  Даны функции
  \[
  f(x)=3x-2,
  \qquad
  g(x)=1-x.
  \]
  Найдите
  \[
  f(g(x))
  \quad\text{и}\quad
  g(f(x)).
  \]
  Определите характер монотонности каждой полученной функции.

  \item
  Представьте функцию
  \[
  y=\left|x^2-4x+1\right|
  \]
  как композицию $f(g(x))$, указав внешнюю и внутреннюю функции.

  \item
  Для функции
  \[
  y=\sqrt{5-2x}
  \]
  найдите область определения и определите характер монотонности.

\end{enumerate}

\textbf{Уровень выше среднего}

\begin{enumerate}[resume]

  \item
  Даны функции
  \[
  f(x)=\sqrt{x},
  \qquad
  g(x)=3-x.
  \]

  Найдите
  \[
  f(g(x))
  \quad\text{и}\quad
  g(f(x)),
  \]
  укажите область определения каждой композиции
  и определите её монотонность.

  \item
  Исследуйте область определения и монотонность функции
  \[
  y=\frac1{x+4}-\sqrt{x+1}.
  \]

  \item
  Исследуйте область определения и монотонность функции
  \[
  y=\frac1{x^5}+\frac1{\sqrt{x}}.
  \]

  \item
  Найдите уравнение прямой, проходящей через точки
  \[
  A(-2,5),
  \qquad
  B(4,-7).
  \]
  Укажите характер её монотонности.

  \item
  Найдите точки пересечения графика
  \[
  y=-3x+6
  \]
  с осями координат.

  \item
  Прямая проходит через точки
  \[
  A(1,-2),
  \qquad
  B(5,6).
  \]
  Найдите её формулу и определите,
  возрастает она или убывает.

\end{enumerate}

\textbf{Сложный уровень}

\begin{enumerate}[resume]

  \item
  Рассмотрите функцию
  \[
  H(x)=\frac1{\sqrt{6-2x}+1}.
  \]

  Представьте её как композицию нескольких простых функций,
  найдите область определения и определите характер монотонности
  на всей области определения.

\end{enumerate}

\clearpage

\sect{6. Ответы к тренировочному блоку}

\begin{center}
\begin{tabularx}{\linewidth}{@{}cX@{}}
\toprule
№ & Ответ \\
\midrule

1
&
\[
f(g(x))=1-3x,
\qquad
g(f(x))=3-3x.
\]
Обе функции убывают.
\\

2
&
Например,
\[
f(t)=|t|,
\qquad
g(x)=x^2-4x+1.
\]
\\

3
&
\[
D=\left(-\infty,\frac52\right].
\]
Функция убывает.
\\

4
&
\[
f(g(x))=\sqrt{3-x},
\qquad
D_{f\circ g}=(-\infty,3],
\]
функция убывает;
\[
g(f(x))=3-\sqrt{x},
\qquad
D_{g\circ f}=[0,+\infty),
\]
функция убывает.
\\

5
&
\[
D=[-1,+\infty).
\]
Функция убывает на всей области определения.
\\

6
&
\[
D=(0,+\infty).
\]
Функция убывает.
\\

7
&
\[
y=-2x+1.
\]
Функция убывает.
\\

8
&
Пересечение с $Ox$:
\[
(2,0).
\]
Пересечение с $Oy$:
\[
(0,6).
\]
\\

9
&
\[
y=2x-4.
\]
Функция возрастает.
\\

10
&
\[
D=(-\infty,3].
\]

Можно взять последовательность:
\[
u(x)=6-2x,
\qquad
v(t)=\sqrt{t},
\qquad
w(s)=s+1,
\qquad
q(r)=\frac1r.
\]

На рассматриваемых промежутках:
$u$ убывает,
$v$ возрастает,
$w$ возрастает,
$q$ убывает.

Итоговая функция $H$ возрастает на
\[
(-\infty,3].
\]
\\

\bottomrule
\end{tabularx}
\end{center}

\end{document}